\section{Ampliación hacia los VLA modernos: de la recipe de 2023 a los modelos unificados}

Esta sección es una interpretación de ingeniería escrita en 2026 y no forma parte de lo que se expuso en la clase de 2023. No usa sistemas posteriores para demostrar retrospectivamente que esta clase tenía razón; analiza qué restricciones de la recipe de Ted Xiao siguen vigentes en un vision-language-action model unificado.

\subsection{Un modelo unificado sigue necesitando una validación modular}

Aunque la visión, el lenguaje y la acción se integren en un solo modelo, el entrenamiento y el despliegue deben seguir auditando por separado perception, planning, skill execution, feedback y data coverage. Una loss de extremo a extremo permite compartir representaciones, pero hace más difícil localizar el origen de los errores; por eso, la evaluation a nivel de sistema debe conservar la verificación de ejecutabilidad al estilo de SayCan, el ciclo cerrado al estilo de Inner Monologue y la auditoría semántica de datos al estilo de DIAL.

\begin{importantbox}{Parámetros unificados no implican responsabilidades unificadas}
Un checkpoint puede cumplir varias funciones a la vez, pero el contrato de ingeniería todavía debe responder: qué clase de token puede controlar acciones reales, quién verifica que una acción sea segura, quién mantiene el estado a largo plazo, quién decide que la tarea terminó y qué entradas son solo datos y no instrucciones. Compartir parámetros no sustituye los límites de permisos ni los de validación.
\end{importantbox}

\subsection{La Data Mixture requiere una Resource Accounting explícita}

La mezcla moderna de datos de varios robots debe registrar la proporción de cada embodiment, tarea, entorno y etiqueta de lenguaje. Si una tarea sencilla y frecuente ocupa la mayor parte de los tokens, la disminución de la loss de entrenamiento puede ocultar el deterioro de habilidades escasas pero críticas; si los distintos action spaces no se normalizan, el modelo aprenderá a asociar un mismo token con controles contradictorios. Los resultados de task-diversity de RT-1 siguen recordándonos que el token count no equivale a cobertura.

\subsection{Memory y Latency son restricciones conjuntas}

Un context más largo permite conservar más historial visual, pero aumenta el costo de la pasada hacia adelante y la latencia de control. Los sistemas corporeizados suelen necesitar organizar en capas el estado reciente de alta frecuencia, los resúmenes de la tarea a largo plazo y una episodic memory recuperable, en lugar de mantener todos los cuadros de forma permanente en la attention window. Al evaluar, deben reportarse a la vez success, latency, control frequency, memory size y recovery behavior.

\begin{warningbox}{Límites de la interpretación actualizada}
Los sistemas de la clase operaban en una cocina cerrada, con acciones limitadas y detectors especializados; un VLA de mundo abierto enfrenta más sensores, la seguridad de las personas y errores irrecuperables. Lo que se transfiere es «la disciplina de interfaces y evidencia», no una extrapolación directa de las tasas de éxito de la clase a despliegues domésticos o industriales.
\end{warningbox}

\subsection{Resumen de la sección}

Un VLA unificado puede absorber varios módulos de la recipe de 2023, pero sigue requiriendo una validación independiente de la ejecutabilidad, la retroalimentación, la cobertura de datos y el presupuesto de tiempo real. Cuanto más sólida sea la representación unificada del Foundation model, con más transparencia debe reportar el sistema qué sigue dependiendo de data, controllers o safety layers especializados.
